\section{A bound on the order of vanishing}
\label{sec:multiplicity}

The preceding representation turns a formula into a scalar component of a
rational differential system. A multiplicity estimate of
Moura~\cite{Moura2003} bounds how long such a component can vanish.
We give a self-contained formal proof of the required estimate,
valid directly over any field of characteristic zero.

For a nonzero series $b\in F[[z]]$, write $\operatorname{ord}_0 b$ for the
least exponent with nonzero coefficient, and put
$\operatorname{ord}_0 0=+\infty$.  For $b_1,\ldots,b_k\in F[[z]]$,
their \emph{Wronskian} is
\[
 W(b_1,\ldots,b_k)
 =\det\bigl(b_\ell^{(j)}\bigr)_{
       0\le j<k,\,1\le\ell\le k}.
\]
The proof of the formal power-series Wronskian criterion by
Bostan and Dumas~\cite[Theorem~2 and Lemmas~1--3]{BostanDumas2010}
uses a basis with distinct orders and a Vandermonde determinant.
We recall that argument, recording also the order of the Wronskian.

\begin{lemma}\label{lem:wronskian}
Let $F$ be a field of characteristic zero and let $V\subset F[[z]]$ be a
vector space of finite positive dimension $k$ over $F$.  There is a basis
$b_1,\ldots,b_k$ of $V$ with distinct orders
$0\le a_1<\cdots<a_k$.  For this basis,
\begin{equation}\label{eq:wronskian-order}
 \operatorname{ord}_0 W(b_1,\ldots,b_k)
 =\sum_{\ell=1}^k a_\ell-\binom{k}{2}.
\end{equation}
In particular, the Wronskian of any basis of $V$ is nonzero, its order
does not depend on the basis, and every nonzero member of $V$ has order
at most $a_k$.
\end{lemma}

\begin{proof}
Consider the decreasing filtration $V_r=V\cap z^rF[[z]]$ for $r\ge0$.
The coefficient of $z^r$ embeds $V_r/V_{r+1}$ into $F$, so each quotient
has dimension at most one.  Moreover, $\bigcap_r V_r=0$.  Since $V$ is
finite-dimensional, the filtration is eventually zero: otherwise its
dimensions would eventually stabilize at a positive value, and then the
subspaces themselves would stabilize, contradicting the zero
intersection.  Choosing a vector at each of the $k$ dimension drops
gives a basis with the asserted distinct orders.  A nonzero linear
combination has the order of its lowest-order basis term.

Write $b_\ell=\beta_\ell z^{a_\ell}+\cdots$, with
$\beta_\ell\ne0$.  Factoring $z^{a_\ell}$ from column $\ell$ and $z^{-j}$
from row $j$ in the Laurent-series field, the remaining constant
determinant is
\[
 \left(\prod_{\ell=1}^k\beta_\ell\right)
 \det\bigl((a_\ell)_j\bigr)_{j,\ell}
 =\left(\prod_{\ell=1}^k\beta_\ell\right)
   \prod_{r<\ell}(a_\ell-a_r).
\]
Here $(a)_0=1$ and $(a)_j=a(a-1)\cdots(a-j+1)$; these are monic
polynomials of successive degrees, so their determinant is the
Vandermonde determinant.  It is nonzero in characteristic zero.  The
factorization remains valid when $a_\ell<j$: the corresponding falling
factorial vanishes and the remaining entry is still a power series.
This proves \eqref{eq:wronskian-order}.  A change of basis over $F$
multiplies the Wronskian by the nonzero determinant of the constant
change-of-basis matrix, proving the remaining assertions.
\end{proof}

The bound below is the specialization of Moura's estimate to the present
degree assumptions. In its complex formulation, recalled in
\cite[Section~2, Theorem~1]{MouraHyperelliptic2003}, the bound is
$m-1+\binom m2(\sigma-2)$, where $\sigma$ is the sum of pole orders
of the matrix differential form $(P/q)\,dz$ on the projective line.
The finite poles contribute at most $n$, and infinity contributes at
most one. The case $P=0$ is immediate, since every solution is constant.
The argument below proves the stated bound formally and follows the
same derivative-row and Wronskian method.

\begin{theorem}[Multiplicity estimate of Moura]\label{thm:multiplicity}
Let $F$ be a field of characteristic zero, and let $n,m\ge1$.
Suppose that $q\in F[z]$ has degree $n$ and $q(0)\ne0$, that
$P\in\operatorname{Mat}_m(F[z])$ has entries of degree at most $n-1$,
and that
\[
 G\in\operatorname{GL}_m(F[[z]]),\qquad G'=(P/q)G.
\]
For a constant row $u\in F^{1\times m}$ and a constant column
$v\in F^{m\times1}$, if $y=uGv$ is nonzero, then
\begin{equation}\label{eq:multiplicity-bound}
 \operatorname{ord}_0 y\le(n-1)\binom{m}{2}+m-1.
\end{equation}
\end{theorem}

\begin{proof}
Put $h=uG$, and let $V$ be the $F$-span of its $m$ component series.
Since $y\ne0$, its dimension $k$ satisfies $1\le k\le m$.
We compare Wronskians of these components with minors of a rational
matrix whose numerator degrees are controlled.

Define rational rows
\[
 R_0=u,\qquad R_{j+1}=R_j'+R_jP/q.
\]
Then $h^{(j)}=R_jG$.  Writing $R_j=Z_j/q^j$, the recurrence becomes
\[
 Z_0=u,\qquad
 Z_{j+1}=qZ_j'-jq'Z_j+Z_jP.
\]
Induction shows that $Z_j$ is a polynomial row with entries of degree
at most $j(n-1)$.  This includes $n=1$, when every $Z_j$ is constant.
Every $R_j$ is regular at zero, because $q(0)\ne0$.

Let $B$ be the $k\times m$ matrix with rows
$R_0,\ldots,R_{k-1}$, and put $S=\binom{k}{2}$.
The rows of $BG$ are $h,h',\ldots,h^{(k-1)}$.
Some $k$ components of $h$ form a basis of $V$, so
Lemma~\ref{lem:wronskian} gives a nonzero $k\times k$ minor of $BG$.
Since $G$ is invertible, $B$ also has a nonzero $k\times k$ minor.
Every such minor has denominator $q^S$ and polynomial numerator of
degree at most $(n-1)S$.  As $q(0)\ne0$, the minimum order $\delta$
among the nonzero $k\times k$ minors of $B$ satisfies
\begin{equation}\label{eq:minor-order}
 0\le\delta\le(n-1)S.
\end{equation}

The minimum order of the $k\times k$ minors of $BG$ is also $\delta$.
Indeed, Cauchy--Binet expresses every such minor as an $F[[z]]$-linear
combination of the minors of $B$, so its order is at least $\delta$.
Applying the same argument to $B=(BG)G^{-1}$ gives the reverse
inequality, because $G^{-1}$ also has power-series entries.
Equivalently, multiplication by $G$ preserves the ideal generated by
these minors in $F[[z]]$.

A nonzero minor of $BG$ attaining this minimum is a Wronskian of $k$
components of $h$.  Those components are linearly independent over
$F$, and hence form a basis of $V$.  Replace them by a basis of distinct
orders $a_1<\cdots<a_k$ as in Lemma~\ref{lem:wronskian}.  Its Wronskian
still has order $\delta$, so
\[
 \delta=\sum_{j=1}^k\bigl(a_j-(j-1)\bigr).
\]
Every summand is nonnegative.  Since $y\in V$ is nonzero, the same
Lemma and \eqref{eq:minor-order} give
\[
 \operatorname{ord}_0 y\le a_k
 \le\delta+k-1
 \le(n-1)\binom{k}{2}+k-1
 \le(n-1)\binom{m}{2}+m-1.
\]
\end{proof}

For the formula systems of the preceding section,
$q(z)=\prod_{i=1}^n(z+i)$ and $G(0)=I$, so every hypothesis of
Theorem~\ref{thm:multiplicity} holds.  Word-series independence makes
$uGv$ nonzero for a nonzero formula polynomial; the theorem bounds the
number of Taylor coefficients needed to detect it.  We now realize
these coefficients by finite rational matrices.
